\documentclass{article}

\textwidth=6.5in
\textheight=8.5in
\voffset=-54pt
\oddsidemargin=0in
\evensidemargin=0in
\setlength{\parskip}{8pt}
\setlength{\parindent}{0pt}

\usepackage{amsmath, amssymb}
\usepackage{ifthen}

\usepackage[inline]{enumitem}
\setlist{topsep=0pt, parsep=8pt}

\usepackage[rgb,table]{xcolor}\convertcolorsUfalse
\usepackage{graphicx}

% change hyperref options
\PassOptionsToPackage{hyphens}{url}
\usepackage[bookmarks,colorlinks,linkcolor=black,citecolor=blue,pdfstartview=FitH,urlcolor=blue]{hyperref}
\hypersetup{pdfpagemode=UseNone}
\usepackage{bookmark}

\usepackage{graphicx}
\usepackage{multicol}
\usepackage{framed}
\usepackage{array}

\usepackage{icomma}

%
\providecommand{\topic}{}
\renewcommand{\topic}[1]{}

\usepackage{cancel}  
\newcommand{\colorcancel}[2]{\renewcommand{\CancelColor}{\color{#1}}\cancel{#2}\renewcommand{\CancelColor}{\color{black}}}
\usepackage{pifont}

\newlist{altenumerate}{enumerate}{3}

\usepackage{soul}

\usepackage{pgfplots}
\pgfplotsset{compat=1.18}  
\pgfplotscreateplotcyclelist{mycolor}{black, blue, red, green!70!black, magenta, purple, cyan, orange }  
\pgfplotsset{
  disabledatascaling,
  cycle list=tikzcolor, 
  axis lines=middle,
  every axis plot/.style={no marks, thick},
  every axis/.style={
    enlarge x limits=false,
    enlarge y limits=auto,
    xlabel style={at={(current axis.right of origin)},anchor=west},
    ylabel style={at={(current axis.above origin)},anchor=south},
    % extra x and y ticks for asymptotes
    extra x tick style={grid=major, grid style={dashed,black}},
    extra y tick style={grid=major, grid style={dashed,black}},
    /pgf/number format/1000 sep={},
  },    
  tick label style = {font=\footnotesize, fill=\currentbgcolor, inner sep=0pt},    
  xticklabel shift=1pt,    
  scaled ticks=false,
  scaled y ticks=false,
  unbounded coords=jump,
  samples=101,
  minor tick length=0,
  scatter/classes={
    open={mark=*,draw=black,fill=white, mark size=2, thin},
    closed={mark=*,draw=black,fill=black, mark size=2, thin}
  },
  width=3in,
}
\pgfplotsset{open/.style={only marks, mark=*, fill=white, thin, mark size=2}}
\pgfplotsset{closed/.style={only marks, mark=*, draw=black, fill=black,
    thin, mark size=2}}
\pgfplotsset{labeled/.style={nodes near coords, only marks, mark=*, draw=black, fill=black, thin, mark size=2, point meta=explicit symbolic}}

\pgfplotsset{gridgraph/.style={clip=false, axis line style={shorten
      >=-8pt}, xlabel style={xshift=8pt}, ylabel style={yshift=8pt},
    enlargelimits=false, grid=major, tick style={black}}} 

\newcommand{\axes}[1][]{
  \begin{tikzpicture}
    \begin{axis}[xtick=\empty, ytick=\empty, ymin=-2, ymax=2,
      xmin=-4.25, xmax=4.25, #1]
    \end{axis}
  \end{tikzpicture}
}
\usepackage{tikz}
\usetikzlibrary{
  matrix,%
  % enables the snake paths
  decorations.pathmorphing,%
  % enables bigger arrows
  decorations.markings,%
  decorations.pathreplacing,%
  decorations.text,%
  calc,%
  patterns,%
  shapes.geometric,%
  arrows,
  decorations.shapes,
  positioning,
  plotmarks,
  intersections
}
\tikzset{>=latex} % sets default arrow type in tikz
\tikzset{font=\normalsize}
\tikzset{mark size=1.5pt, mark options=thin}
\tikzset{pin distance=4pt,
  every pin edge/.style={<-, thin, shorten <= -2pt}}

\interfootnotelinepenalty=10000
\renewcommand{\thefootnote}{(\arabic{footnote})}

\usepackage{multirow, bigdelim}

% change to \usepackage[sols]{handout} to see solutions
\usepackage[sols]{handout}

\newcommand\iscurrentenumi[1]{%
  \ifthenelse{\equal{#1}{\theenumi}}{}{#1}}
\setenumerate[1]{ref=\#\arabic*}
\setenumerate[2]{ref=\protect\iscurrentenumi{\theenumi}(\alph*)}
\newcommand\enumiiprefix[2]{%
  \ifthenelse{{\equal{#1}{\theenumi}}\AND{\equal{#2}{(\alph{enumii})}}}{}{}%
  \ifthenelse{{\equal{#1}{\theenumi}}\AND{\NOT{\equal{#2}{(\alph{enumii})}}}}{#2}{}%
  \ifthenelse{\equal{#1}{\theenumi}}{}{#1#2}%
}
\setenumerate[3]{ref=\protect\enumiiprefix{\theenumi}{(\alph{enumii})}\roman*}
% Make an environment that puts items in columns, first going across. 
\usepackage{tabto}
\newenvironment{enumeratecols}[2][]
{\NumTabs{#2}%
  \begin{enumerate*}[%before={\hspace{\dimexpr-\parindent-1pt}\tab},%
    itemjoin={\tab},#1]}%
  {\end{enumerate*}}

\makeatletter

% underline links               
\let\@href=\href
\renewcommand{\href}[2]{\@href{#1}{\ul{#2}}}

\newdimen\SOUL@dimen %new
\def\SOUL@ulunderline#1{{%
    \setbox\z@\hbox{#1}%
    % \dimen@=\wd\z@
    \SOUL@dimen=\wd\z@ %new
    \dimen@i=\SOUL@uloverlap
    \advance\SOUL@dimen2\dimen@i %\dimen@ exchanged too
    \rlap{%
      \null
      \kern-\dimen@i
      % \SOUL@ulcolor{\SOUL@ulleaders\hskip\dimen@}%
      \SOUL@ulcolor{\SOUL@ulleaders\hskip\SOUL@dimen}% new
    }%
    \unhcopy\z@
  }}

\let\@uhref=\href
\usepackage[style=verbose-ibid,backend=bibtex]{biblatex}
\let\@autocite=\autocite
\renewcommand{\autocite}[1]{
  \let\href=\@href
  \@autocite{#1}
  \let\href=\@uhref
}
\let\@autocites=\autocites
\renewcommand{\autocites}[2]{
  \let\href=\@href
  \@autocites{#1}{#2}
  \let\href=\@uhref
}
\bibliography{references}
\makeatother

\renewcommand{\answer}[1]{#1}

\definecolor{purple}{rgb}{0.5,0,1.0}
\definecolor{hotpink}{rgb}{1, 0, 0.5}

\definecolor{skyblue}{rgb}{0, 0.5, 1}

\newcommand{\highlight}[1]{\boldmath\hl{\textbf{#1}}\unboldmath}

\newcommand{\goal}[1]{\textbf{\textcolor{skyblue}{#1}}}

\usepackage{amsthm}

% make URLs print in the same font
\usepackage{url}
\urlstyle{same}

\usepackage{pifont}
\def\exend{\hfill\ding{118}}

%\makeatletter\@addtoreset{equation}{section}\makeatother
%\renewcommand{\theequation}{\arabic{section}.\arabic{equation}}

\newtheoremstyle{theorem}% name
{8pt}%      Space above
{0pt}%      Space below
{\itshape}%         Body font
{}%         Indent amount (empty = no indent, \parindent = para indent)
{\bfseries}% Thm head font
{.}%        Punctuation after thm head
{.5em}%     Space after thm head: " " = normal interword space;
                                %       \newline = linebreak
{}%         Thm head spec (can be left empty, meaning `normal')

\theoremstyle{theorem}

\let\Theorem\undefined
\newtheorem{Theorem}[equation]{Theorem}
\let\Definition\undefined
\newtheorem{Definition}[equation]{Definition}
\let\Summary\undefined
\newtheorem{Summary}[equation]{Summary}
\let\Conjecture\undefined
\newtheorem{Conjecture}[equation]{Conjecture}
\let\Fact\undefined
\newtheorem{Fact}[equation]{Fact}

\renewcommand{\thesubsection}{\arabic{subsection}}
\makeatletter
\def\@seccntformat#1{\@ifundefined{#1@cntformat}%
   {\csname the#1\endcsname\quad}%       default
   {\csname #1@cntformat\endcsname}}%    enable individual control
\newcommand\section@cntformat{}
\makeatother

  \usepackage{exam}

\usepackage{fancyhdr}
\lhead{\textcolor{white}{This content is protected and may not be shared, uploaded, or distributed.}}
\rhead{}
\rfoot{\vspace*{\fill}\footnotesize\textcopyright{} \emph{Harvard Math Ma}}
\renewcommand{\headrulewidth}{0pt}
\pagestyle{fancy}
\begin{document}

  \title{Exam 2, Practice 3}

\instructions{instructions.tex}{}

\listofproblems

\begin{enumerate}
\item \points{12} 

  Your friend would like to approximate $\sqrt{8.4}$ by hand and thinks of 3 ways to do this:
  \begin{altenumerate}[label=\arabic*.]
  \item Use the secant line of $y = \sqrt{x}$ between $x = 4$ and $x = 9$.
  \item Use the secant line of $y = \sqrt{x}$ between $x = 9$ and $x = 16$.
  \item Use the line tangent to $y = \sqrt{x}$ at $x = 9$.
  \end{altenumerate}

  \begin{enumerate}
  \item
    \begin{problem}[\vfill]
      Sketch 3 graphs of $\sqrt{x}$; on the three graphs, illustrate your friend's 3 methods.
      In addition, say whether each method will give an upper bound or
      lower bound for $\sqrt{8.4}$.
    \end{problem}

    \begin{solution}
      In each picture, the approximation of $\sqrt{8.4}$ is the height of the line at $x = 8.4$.
      \begin{center}
        \begin{tikzpicture}[declare function={f(\x)=ln(1+\x);}]
          \begin{axis}[xtick={4, 9, 16}, width=2.5in, ytick=\empty]
            \addplot+[domain=0:18]{ln(1+x)}; 
            \addplot+[sharp plot, hotpink] coordinates {(4, {f(4)}) (9, {f(9)})} node[pos=0.5, anchor=north west]{Method 1}; %
          \end{axis}
        \end{tikzpicture}
        \qquad
        \begin{tikzpicture}[declare function={f(\x)=ln(1+\x);}]
          \begin{axis}[xtick={4, 9, 16}, width=2.5in, ytick=\empty]
            \addplot+[domain=0:18]{ln(1+x)}; 
            \addplot+[domain=7:18, green!80!black] {(f(16)-f(9))/7*(x-9) + f(9)} node[pos=0, anchor=south east]{Method 2}; %
          \end{axis}
        \end{tikzpicture}
        \qquad
        \begin{tikzpicture}[declare function={f(\x)=ln(1+\x);}]
          \begin{axis}[xtick={4, 9, 16}, width=2.5in, ytick=\empty]
            \addplot+[domain=0:18]{ln(1+x)}; 
            \addplot+[domain=6:12, skyblue] {1/10*(x-9) + f(9)} node[pos=1, anchor=south east]{Method 3}; % 
          \end{axis}
        \end{tikzpicture}        
      \end{center}
      We can see that \fbox{Method 1 will give a lower bound} and \fbox{Methods 2 and 3 will give upper bounds}. 
    \end{solution}

  \item
    \begin{problem}[\nstretch{0.5}]
      Based on your pictures, which method do you expect to give the
      closest approximation? 
    \end{problem}

    \begin{solution}
      It looks like \answer{Method 3} will be closest. 
    \end{solution}

  \item
    \begin{problem}[\vfill]
      Calculate the approximation your friend's 3rd idea gives. 
    \end{problem}

    \begin{solution}
      We need to find the tangent line at $x = 9$. The slope is $f'(9)$:
      \begin{align*}
        f'(x) &= \dfrac{1}{2} x^{-1/2} = \dfrac{1}{2 \sqrt{x}} \\
        f'(9) &= \dfrac{1}{2 \sqrt{9}} = \dfrac{1}{6}
      \end{align*}
      Since $(9, 3)$ is a point on the tangent line, the equation of the line is $y - 3 = \dfrac{1}{6} (x - 9)$, or $y = 3 + \dfrac{1}{6} (x - 9)$.
      So,
      \begin{align*}
        \sqrt{x} & \approx 3 + \dfrac{1}{6} (x - 9) \text{ for $x$ near 9} \\
        \intertext{Plugging in $x = 8.4$ gives}
        \sqrt{8.4} & \approx 3 + \dfrac{1}{6} (8.4 - 9) \\
        & = 3 - 0.1 \\
        &= \boxed{2.9}
      \end{align*}
    \end{solution}

  \item
    \begin{problem}[\vfill\newpage]
      Another friend suggests that you could get an even better approximation by hand if you use the secant line of $y = \sqrt{x}$ between $x = 8.3$ and $x = 8.5$. Explain why this is or isn't a good idea.
    \end{problem}

    \begin{solution}
      \answer{This approximation would involve $\sqrt{8.3}$ and $\sqrt{8.5}$; we don't know the exact values of those, so we wouldn't get a very useful approximation!}
    \end{solution}

  \end{enumerate}

\item \points{14}

  \begin{problem}[\newpage]
    A company called The Pi(e) Guys is designing a metal can for pie filling. The metal used for the sides costs 1 cent per cm$^2$, while the metal used for the top and bottom costs 2 cents per cm$^2$. Since the founders love both pie and $\pi$, they have decided to spend $300 \pi$ cents on each can. If they want to maximize the volume of each can, what dimensions should they make the can?

    Please communicate your reasoning clearly, and organize your work so the reader can follow it easily. 
  \end{problem}

  \begin{solution}
    \highlight{Goal:} Maximize \goal{volume}

    \highlight{Define variables:}
    \begin{itemize}[nosep]
    \item \goal{$V$ = volume}
    \item $r$ = radius
    \item $h$ = height
    \end{itemize}
    Here are two possible cans:
    \begin{center}
      \begin{tikzpicture}[x=0.75cm, y=0.75cm]
        \draw (0,0) ellipse (1.25 and 0.5);
        \draw[dashed] (0, 0) -- node[midway, above, inner sep=1pt] {$r$} (1.25, 0); 
        \draw (-1.25,0) -- (-1.25,-3.5);
        \draw (-1.25,-3.5) arc (180:360:1.25 and 0.5);
        \draw [dashed] (-1.25,-3.5) arc (180:360:1.25 and -0.5);
        \draw (1.25,-3.5) -- (1.25,0);
        \draw[|-|] (-1.75, -3.5) -- node[midway, inner sep=6pt, fill=white] {$h$} (-1.75,0); 
      \end{tikzpicture}
      \hspace{0.25in}
      \begin{tikzpicture}[x=1.25cm, y=0.375cm]
        \draw (0,0) ellipse (1.25 and 0.5);
        \draw[dashed] (0, -3.5) -- node[midway, above, inner sep=1pt] {$r$} (1.25, -3.5); 
        \draw (-1.25,0) -- (-1.25,-3.5);
        \draw (-1.25,-3.5) arc (180:360:1.25 and 0.5);
        \draw [dashed] (-1.25,-3.5) arc (180:360:1.25 and -0.5);
        \draw (1.25,-3.5) -- (1.25,0);
        \draw[|-|, xshift=10pt] (-1.75, -3.5) -- node[midway, fill=white] {$h$} (-1.75,0); 
      \end{tikzpicture}
    \end{center}        
    \highlight{Formula for goal:} $\goal{$V$} = \pi r^2 h$

    This formula involves both $r$ and $h$, so we need to \highlight{relate $r$ and $h$ using something else we're told in the problem}.
    \begin{align*}
      \text{cost} &= 300 \pi \\
      \text{cost of top and bottom} + \text{cost of sides} &= 300 \pi \\
      (2 \text{ cents / cm$^2$})(\text{surface area of top and bottom}) + (1 \text{ cents/cm$^2$})(\text{surface area of sides}) &= 300 \pi \\
      2 (2 \pi r^2) + 1 (2 \pi r h) &= 300 \pi \\
      4 \pi r^2 + 2 \pi r h &= 300 \pi
    \end{align*}
    It looks easier to solve this for $h$:
    \begin{align}
      2 \pi r h &= 300 \pi - 4 \pi r^2 \nonumber \\
      h &= \frac{300 \pi - 4 \pi r^2}{2 \pi r} \nonumber \\
      &= \frac{\colorcancel{hotpink}{2 \pi} (150 - 2 r^2)}{\colorcancel{hotpink}{2 \pi} r} \nonumber \\
      &= \frac{150 - 2r^2}{r} \label{eq:fall2023-e2practice-can-optimization}
    \end{align}
    Plugging this into out formula for \goal{$V$}, we get \goal{$V$} as a function of just $r$:
    \[ \goal{$V(r)$} = \pi r^2 \cdot \frac{150 - 2r^2}{r} = \pi r (150 - 2r^2) = \pi(150 r - 2r^3).\]
    \highlight{Critical points:} $V'(r) = \pi (150 - 6r^2)$. This is never undefined, and it's 0 when
    \begin{align*}
      150 &= 6r^2 \\
      25 &= r^2 \\
      \pm 5 &= r
    \end{align*}
    Of course, $r = -5$ doesn't make sense in our situation, so $r = 5$ is the only critical point in the domain. 

    \highlight{Global maximum:} Since there's only one critical point in the domain, we can try the First or Second Derivative Test. Let's try the Second Derivative Test, since it's easy to calculate that $V''(r) = \pi (-12r)$. Then, $V''(5) = \pi (-60) < 0$, so $V(r)$ has a local maximum at $r = 5$. Since this was the only critical point in the domain, this must be the global maximum.

    \highlight{Answer the question:} We can plug $r = 5$ back into (\ref{eq:fall2023-e2practice-can-optimization}) to find that the corresponding value of $h$ is $h = 20$. So, \fbox{the can should have radius 5 cm and height 20 cm}.
  \end{solution}

\item \points{11}

  \begin{enumerate}
  \item
    \begin{problem}
      On the axes below, sketch the graph of $f(x) = 6 \left( \frac{2}{3} \right)^x$. Label any intercepts. In addition, label 2 points with their (simplified) exact coordinates: please label one point with $x > 0$ and one with $x < 0$.
    \end{problem}

    \begin{problemonly}
      \axes[width=4.5in, height=4.5in]
    \end{problemonly}

    \begin{solution}
      Since $\frac{2}{3} < 1$, this is an exponential decay function. The $y$-intercept is $f(0) = 6 \left( \frac{2}{3} \right)^0 = 6$.  To get two points, we can just plug in a few values of $x$. For example, $f(1) = 6 \left( \frac{2}{3} \right)^1 = 6 \cdot \frac{2}{3} = 4$ and $f(-1) = 6 \left( \frac{2}{3} \right)^{-1} = 6 \cdot \frac{3}{2} = 9$. 
      \begin{center}
        \answer{
          \begin{tikzpicture}
            \begin{axis}[xtick=\empty, ytick=\empty, ymax=20.25, ymin=-20.25, width=4.5in, height=4.5in]
              \begin{scope}[skyblue]
                \addplot+[domain=-3:3] {6 * (2/3)^x}; %
                \addplot+[closed] coordinates {(0, 6)} node[anchor=south west]{$(0, 6)$}; %
                \addplot+[closed] coordinates {(1, 4)} node[anchor=north east]{$(1, 4)$}; %
                \addplot+[closed] coordinates {(-1, 9)} node[anchor=south west]{$(-1, 9)$}; %
              \end{scope}
            \end{axis}
          \end{tikzpicture}
        }
      \end{center}
    \end{solution}

  \item
    \begin{problem}[\vfill]
      Using a dashed line, sketch the graph of $g(x) = 6 \left( \frac{2}{3} \right)^{x / 2}$ on the axes in (a). Label any intercepts, and label 2 additional points with their (simplified) exact coordinates. Explain how this graph relates to the graph you drew in (a). 
    \end{problem}

    \begin{solution}
      Notice that $g(x) = f\left( \dfrac{x}{2} \right)$, so we get the graph of $g$ by stretching the graph of $f$ horizontally by a factor of $2$. This means the point $(0, 6)$ stays where it is, the point $(1, 4)$ moves to $\left( 2, 4 \right)$, and the point $(-1, 9)$ moves to $\left( -2, 9 \right)$:
      \begin{center}
        \answer{
          \begin{tikzpicture}
            \begin{axis}[xtick=\empty, ytick=\empty, ymax=20.25, ymin=-20.25, width=4.5in, height=4.5in]
              \begin{scope}[skyblue]
                \addplot+[domain=-3:3] {6 * (2/3)^x}; %
                \addplot+[closed] coordinates {(1, 4)} node[anchor=north east]{$(1, 4)$}; %
                \addplot+[closed] coordinates {(-1, 9)} node[anchor=south west]{$(-1, 9)$}; %
              \end{scope}
              \addplot+[dashed, domain=-3:3] {6 * (2/3)^(x / 2)}; % 
              \addplot+[closed] coordinates {(0, 6)} node[anchor=south west]{$(0, 6)$}; %
              \addplot+[closed] coordinates {(2, 4)} node[anchor=south west, fill=white]{$\left(2, 4 \right)$}; %
              \addplot+[closed] coordinates {(-2, 9)} node[anchor=north east, fill=white]{$\left(-2, 9 \right)$}; %
            \end{axis}
          \end{tikzpicture}
        }
      \end{center}
    \end{solution}

  \item Suppose $a$ and $b$ are numbers such that $3^a = 2$ and $3^b = 5$. Simplify the following expressions. (Your answers shouldn't involve $a$ and $b$.) 

    \begin{enumerate}
    \item
      \begin{problem}[\vfill]
        $3^{b - a}$
      \end{problem}

      \begin{solution}
        $3^{b - a} = \dfrac{3^b}{3^a} = \boxed{\frac{5}{2}}$.
      \end{solution}

    \item
      \begin{problem}[\vfill\newpage]
        $3^{a + 2b}$
      \end{problem}

      \begin{solution}
        \begin{align*}
          3^{a + 2b} &= 3^a \cdot 3^{2b} \\
          &= 3^a \cdot (3^b)^2 \\
          &= 2 \cdot 5^2 \\
          &= \boxed{50}
        \end{align*}
      \end{solution}

    \end{enumerate}

  \end{enumerate}

\item \label{fall2014-1a-midterm1-5-modified} \points{15}  % Jason Bland

  The graph of a function $f$ defined on $[-3, 3]$ is given below:
  \begin{center}
    \begin{tikzpicture}
      \begin{axis}[gridgraph, xtick={-3, -2, ..., 3}]
        \addplot+[domain=-3:-2] {3 + x};%
        \addplot+[domain=-2:0] {2 + 2 * x};%
        \addplot+[domain=0:1] {-1- x};%
        \addplot+[domain=1:3] {-2 * (x - 2)^2};%
        \addplot[closed] coordinates { (-3, 0) (-2, -2) (0, -1) (3, -2) };%
        \addplot[open] coordinates { (-2, 1) (0, 2) (-0.5, 1)};%
      \end{axis}
    \end{tikzpicture}
  \end{center}

  \begin{enumerate}
  \item 
    \begin{problem}[\stretch{2}]
      Sketch a graph of $f'(x)$ on $(-3, 3)$. Make sure it's clear where $f'$ is undefined. 
    \end{problem}

    \begin{solution}
      $f$ is not differentiable at $x = -2, -0.5, 0, 1$, so $f'$ will be undefined there. Between $x = -3$ and $x = -2$, the graph of $f$ has slope 1, so the graph of $f'$ should just equal 1 on the interval $(-3, -2)$. Between $x = -2$ and $x = 0$ (except at $x = -0.5$), the graph of $f$ has slope 2, so the graph of $f'$ should equal 2 there. Between $x = 0$ and $x = 1$, the graph of $f$ has slope $-1$, so the graph of $f'$ should equal $-1$ there. Right after $x = 1$, the slope is positive; as $x$ increases, the slope gets less positive, hitting 0 at $x = 2$. After that, it gets more and more negative. So, that gives us a graph like this: 
      \begin{center}
        \answer{
          \begin{tikzpicture}
            \begin{axis}[xtick={-3, -2, ..., 3}, ytick={-1, 1, 2}]
              \addplot+[domain=-3:-2] {1};%
              \addplot+[domain=-2:0] {2};%
              \addplot+[domain=0:1] {-1};%
              \addplot+[domain=1:3] {-4 * (x - 2)};%
              \addplot+[open] coordinates { (-3, 1) (-2, 1) (-2, 2) (-0.5, 2) (0, 2) (0, -1) (1, -1) (1, 4) (3, -4) };% 
            \end{axis}
          \end{tikzpicture}
        }
      \end{center}
      If you're having trouble understanding this graph, try playing around with \href{https://www.geogebra.org/m/jbhsbuhb}{this applet}; drag the green dot to see how the slope changes. 
    \end{solution}

    \begin{grading}
      The piece on the interval $(1, 3)$ should be decreasing and have an $x$-intercept of 2, but we don't have enough information to tell whether it's actually linear. 
    \end{grading}

  \item Evaluate the following limits or state that they do not exist. If a limit does not exist but is $\infty$ or $-\infty$, indicate that as well. Show your work / explain your reasoning. 

    \begin{enumerate}

    \item 
      \begin{problem}[\vfill\newpage]
        $\displaystyle \lim_{t \to 2} \frac{3^{-t}}{f(t)}$
      \end{problem}

      \begin{solution}
        As $t \to 2$, the numerator $\to 3^{-2} = \dfrac{1}{9}$ while the denominator $f(t)$ approaches 0. In this situation, we must think about the sign of the denominator.

        When $t$ is just a bit bigger than $2$, the denominator $f(t)$ is just a little bit less than 0 (while the numerator is close to $\frac{1}{9}$), so $\dfrac{3^{-t}}{f(t)}$ will be very negative. Therefore, $\displaystyle \lim_{t \to 2^+} \frac{3^{-t}}{f(t)} = - \infty$.

        When $t$ is just a bit less than $2$, the denominator $f(t)$ is again just a little bit less than 0 (and the denominator is still close to $\frac{1}{9}$), so $\dfrac{3^{-t}}{f(t)}$ will again be very negative. Therefore, $\displaystyle \lim_{t \to 2^-} \frac{3^{-t}}{f(t)} = - \infty$.

        Therefore, $\displaystyle \lim_{t \to 2} \frac{3^{-t}}{f(t)} = \boxed{-\infty}$.
      \end{solution}

    \item \label{fall2014-1a-midterm1-5b}
      \begin{problem}[\vfill]
        $\displaystyle \lim_{x \to -1} \frac{f(x) - f(-1)}{x + 1}$
      \end{problem}

      \begin{solution}
        You should recognize this as $f'(-1)$, which is $\boxed{2}$ from the graph.

        \emph{Alternate approach:} Even if you didn't immediately recognize this as the limit definition of the derivative, you can still do this problem! The expression \( \dfrac{f(x) - f(-1)}{x + 1}\) should jump out as being an average rate of change / slope of a secant line. In particular, it's the slope of the secant line between \((-1, f(-1))\) and $(x, f(x))$. When $x$ is close to $-1$ (which is what we care about since we're taking the limit as $x \to -1$), this slope is 2. So, the limit is also 2. 
      \end{solution}

    \item
      \begin{problem}[\vfill]
        $\displaystyle \lim_{h \to 0^+} \frac{f(-2 + h) - f(-2)}{h}$
      \end{problem}

      \begin{solution}
        This should remind you of the definition of $f'(-2)$, except that this is only a right-handed limit. We can see from the graph that $f'(-2)$ does not exist, but this one-sided limit could still exist.

        You should recognize $\dfrac{f(-2 + h) - f(-2)}{h}$ as the slope of the secant line between the points $(-2, f(-2))$ and $(-2 + h, f(-2 + h))$; when $h > 0$, this means the slope of the secant line between $(-2, f(-2))$ and a point just to the right of that. From the graph, we can see that such a slope is always $2$, so this limit must be \fbox{2} as well.

      \end{solution}

    \item
      \begin{problem}[\vfill]
        $\displaystyle \lim_{w \to 2} \frac{w^2 - 4}{w^2 + 5w + 6}$
      \end{problem}

      \begin{solution}
        As $w \to 2$, the numerator approaches $2^2 - 4 = 0$ and the denominator approaches $2^2 + 5 \cdot 2 + 6 = 20$. So, $\displaystyle \lim_{w \to 2} \frac{w^2 - 4}{w^2 + 5w + 6} = \frac{0}{20} = \boxed{0}$. 
      \end{solution}

    \item
      \begin{problem}[\vfill\newpage]
        $\displaystyle \lim_{w \to -2} \frac{w^2 - 4}{w^2 + 5w + 6}$
      \end{problem}

      \begin{solution}
        As $w \to -2$, the numerator approaches $(-2)^2 - 4 = 0$ and the denominator approaches $(-2)^2 + 5(-2) + 6 = 0$, so we need to factor and cancel.
        \begin{align*}
          \lim_{w \to -2} \frac{w^2 - 4}{w^2 + 5w + 6} &= \lim_{w \to -2} \frac{(w - 2)(w + 2)}{(w + 2)(w + 3)} \\
          &= \lim_{w \to -2} \frac{w - 2}{w + 3} \\
          &= \frac{-4}{1} \\
          &= \boxed{-4}
        \end{align*}
      \end{solution}
    \end{enumerate}

  \end{enumerate}

\item \label{fall2015-1a-m2practice3-1} \label{timolol} \points{16} 

  \topic{graphing}

  The drug timolol treats high blood pressure by making a patient's heart beat more slowly. To prescribe it safely, doctors need to understand how it affects a patient's heart rate while exercising. Based on data, an appropriate model is $\displaystyle E(C) = 160 - \frac{90C}{12 + C}$ where $E$ is exercise heart rate (in beats per minute) and $C$ is the plasma concentration\probonly{\footnote{Plasma concentration basically measures how much drug a patient has in his/her body.}} of timolol (in ng/mL)\solonly{\autocites{RefWorks:12}{Holford1982143}}. We're only interested in this function for $C \geq 0$, since it doesn't make sense to talk about negative drug concentrations.

  \begin{enumerate}
  \item
    \begin{problem}[\vfill]
      On what intervals in $[0, \infty)$ is $E(C)$ increasing? On what intervals is it decreasing?
    \end{problem}

    \begin{solution}
      As usual, we figure this out by looking at the sign of $E'(C)$. Using the Sum Rule and the Quotient Rule,
      \begin{align*}
        E'(C) &= 0 - \frac{(12 + C) 90 - 90C (1)}{(12 + C)^2} \\
        &= - \frac{1080}{(12 + C)^2}
      \end{align*}
      On $[0, \infty)$, this is never undefined or 0, so $E'(C)$ never changes sign. So, we can test any point in $[0, \infty)$ to determine the sign of $E'$ for all of $[0, \infty)$. For example, $E'(1) = - \dfrac{1080}{13^2} < 0$, so $E'(C)$ is negative on $[0, \infty)$.
      \begin{center}
        \begin{tikzpicture}
          \begin{axis}[axis y line=none, x axis line style={-}, xtick={0}, ymin=-2, ymax=2, height=1.25in, width=3in, clip=false]
            \addplot+[domain=-0.1:0.1, very thin]{0};
            \draw(-0.1, 0) node[anchor=south west] {$E$};%
            \draw (-0.1, -1.5) node[right] {sign of $E'$};%
            \draw (0.05, -1.5) node{$-$};%
            \draw[->, xshift=2pt] (0.00666667, 1.5) -- (0.0933333, 0.5);%
          \end{axis}
        \end{tikzpicture}
      \end{center}
      So, $E(C)$ is \fbox{decreasing on $[0, \infty)$}.
    \end{solution}

  \item
    \begin{problem}[\vfill\newpage]
      On what intervals in $[0, \infty)$ is $E(C)$ concave up? On what intervals is it concave down?
    \end{problem}

    \begin{solution}
      Now, we need to calculate $E''(C)$. We need to use the Quotient Rule, but we need to first expand the denominator (since that's the only way we'll be able to differentiate it): $E'(C) = \dfrac{- 1080}{C^2 + 24 C + 144}$. Then, 
      \begin{align*}
        E''(C) &= \dfrac{(C^2 + 24 C + 144)(0) - (-1080) (2C + 24)}{(C^2 + 24 C + 144)^2} \\
        &= \dfrac{1080 \cdot 2 (C + 12)}{(C + 12)^4} \\
        &= \dfrac{2160}{(C + 12)^3}
      \end{align*}
      On $[0, \infty)$, this is never 0 or undefined, so $E''(C)$ never changes sign. So, we can test any point in $[0, \infty)$ to determine the sign of $E''$ everywhere in $[0, \infty)$. For example, $E''(1) = \dfrac{2160}{13^3} > 0$, so $E''(C)$ is always positive. Therefore, $E(C)$ is \fbox{concave up on $[0, \infty)$}.
    \end{solution}

  \item \label{timolol-asymptote}
    \begin{problem}[\stretch{0.5}]
      On $[0, \infty)$, does $E(C)$ have a horizontal asymptote? If so, where? Explain how you know.
    \end{problem}

    \begin{solution}
      We need to calculate $\displaystyle \lim_{C \to \infty} E(C)$, or $\displaystyle \lim_{C \to \infty} \left( 160 - \frac{90C}{12 + C} \right)$. When $C$ is really big, 12 is quite insignificant compared to $C$, so $C + 12 \approx C$, and \[ \lim_{C \to \infty} \left( 160 - \frac{90C}{12 + C} \right) = \lim_{C \to \infty} \left( 160 - \frac{90C}{C} \right) = \lim_{C \to \infty} (160 - 90) = 70.\] Therefore, $E(C)$ has a horizontal asymptote at \fbox{$E = 70$}.
    \end{solution}

  \item
    \begin{problem}[\vfill]
      Sketch the graph of $E(C)$ on $[0, \infty)$. Label all asymptotes and intercepts. 
    \end{problem}

    \begin{solution}
      The $y$-intercept is $E(0) = 160$. We've found that $E(C)$ is decreasing and concave up everywhere on $[0, \infty)$, and $\displaystyle
      \lim_{C \to \infty} E(C) = 70$. So, the graph looks like this:
      \begin{center}
        \answer{
          \begin{tikzpicture}
            \begin{axis}[xtick=\empty, ytick={70, 160}, ymin=0]
              \addplot+[domain=0:200] {160-90*x/(12+x)};
              \addplot+[domain=0:200, dashed, black] {70};
            \end{axis}
          \end{tikzpicture}
        }
      \end{center}
    \end{solution}

  \item
    \begin{problem}[\nstretch{0.5}]
      Explain the practical significance of your answer to
      \ref{timolol-asymptote} in terms of the drug's effect.
    \end{problem}

    \begin{solution}
      \answer{We can see from the graph that, no matter how much timolol the patient takes, their exercise heart rate will never go below 70.}
    \end{solution}

  \end{enumerate}

\item \points{12}  \finishexam

   A smartphone battery, when fully charged, has a maximum voltage of about 4 volts. We can model the voltage of a smartphone battery by the function $V(t) = 4(1-e^{-t})$, where $t$ is the number of hours since the phone started charging. 

\begin{enumerate}
    \item 
    \begin{problem}[\vfill] Calculate $V(0)$. Interpret this value in the context of the problem.
    \end{problem}
\begin{solution} 
$$V(0)=4(1-e^{-0})=4(1-1)=0$$
In the context of the problem, this means that at $t=0$, right when the smartphone begins charging, the voltage of the battery is 0 volts.
\end{solution}
    \item 
    \begin{problem}[\vfill] Determine $\lim\limits_{t\to\infty} V(t)$. Interpret this limit in the context of the problem.\\
    \end{problem}

\begin{solution}
As $t\to\infty$, $e^{-t}\to 0$, so $\displaystyle \lim_{t\to\infty}V(t) = 4$. In the context of the problem, this means that as time continues passing, the battery's voltage approaches 4 volts.
\end{solution}

    \item 
    \begin{problem}[\vfill] Find $V'(t)$.
    \end{problem}
\begin{solution}
    We begin by expanding before differentiating:
        $$V(t) = 4-4e^{-t}$$
        $$V'(t) = 4e^{-t}$$
    \end{solution}
    \item 
    \begin{problem}[\vfill] Calculate $V'(0)$. Interpret this value in the context of the problem.
    \end{problem}
    \begin{solution}
        $$V'(0)=4e^{-0}=4$$
        In the context of this problem, this means at $t=0$ (when the battery begins charging), the voltage is increasing at a rate of 4 volts/hour.
    \end{solution}
    
\begin{problemonly}
\newpage
\end{problemonly}
    \item 
    \begin{problem}[\vfill]
    Determine $\lim\limits_{t\to\infty} V'(t)$. Interpret this limit in the context of the problem.\\
    \end{problem}
\begin{solution}
    Since as $t\to\infty$, $e^{-t}\to 0$, we have:
    $$\lim_{t\to\infty} 4e^{-t}=0$$
    In the context of this problem, this means that as time passes, the rate at which the battery's voltage is changing goes to 0 volts/hour.
\end{solution}
    \item 
    \begin{problem}[\vfill]
    Determine the open intervals on which $V(t)$ is increasing. Justify your answer.\\
    \end{problem}
\begin{solution}
    The intervals on which $V(t)$ is increasing are the intervals on which $V'(t)>0$. \\
    $V'(t) = 4e^{-t}$, which is positive everywhere, \\
    so $V(t)$ is increasing everywhere on the interval $(-\infty, \infty)$.
\end{solution}
    \item 
    \begin{problem}[\vfill] Does $V(t)$ attain a global maximum on the interval $[0, \infty)$? Explain your reasoning.
    \end{problem}

    \begin{solution}
        Since the function is increasing everywhere, it will never achieve a global maxima on the interval $[0,\infty)$.
    \end{solution}

    \item 
    \begin{problem}Sketch the graph of the phone battery's voltage over time. Label all intercepts and asymptotes.
    \end{problem}
     \begin{problemonly}
        \axes[]
      \end{problemonly} 
      \begin{solution}
      This is a function that begins at (0,0) with a slope at 4, continues increasing, approaches 4 in the limit and the slope approaches 0 (meaning the function flattens out).\\
      
          \begin{tikzpicture}
            \begin{axis}[ yticklabels={$4$}, xtick=\empty, ylabel={voltage (volts)}, xlabel={$t$ (hours)}, ]
              \addplot+[domain=0:6] {4-4*e^(-x))};
              
              \addplot[domain=-1:5.5, dashed, red] (\x,4) node[above] {$y=4$};
              \node[above left, outer sep=2pt] (0,0) {(0,0)};
              
            \end{axis}
          \end{tikzpicture}
      \end{solution}

\end{enumerate}
\end{enumerate}

\end{document}
